\section{Mobile Deployment}
\label{sec:deployment}

\subsection{Export Pipeline}

The trained PicoDet-S L1 model is exported and deployed through a
four-stage pipeline:

\begin{figure}[H]
\centering
\resizebox{\textwidth}{!}{%
\begin{tikzpicture}[
  node distance=0.8cm and 0.3cm,
  box/.style={draw, rounded corners=4pt, minimum width=3.0cm,
              minimum height=1.1cm, align=center, font=\small},
  arrow/.style={->, thick, >=Stealth},
  label/.style={font=\scriptsize, gray!80!black}
]
% Stage 1
\node[box, fill=picoBlue!20] (train)
  {\textbf{Training}\\PaddleDetection 2.6\\Tesla T4 $\times$2};

% Stage 2
\node[box, fill=picoBlue!10, right=1.4cm of train] (export)
  {\textbf{Export}\\Paddle Inference\\(w/ \& w/o NMS)};

% Stage 3
\node[box, fill=accentGreen!15, right=1.4cm of export] (onnx)
  {\textbf{ONNX Convert}\\paddle2onnx\\Opset 11};

% Stage 4
\node[box, fill=yoloOrange!15, right=1.4cm of onnx] (ios)
  {\textbf{iOS App}\\ONNX Runtime\\+ CoreML ANE};

% Arrows with labels
\draw[arrow] (train) -- (export)
  node[above, label, midway] {best\_model.pdparams};
\draw[arrow] (export) -- (onnx)
  node[above, label, midway] {pdmodel};
\draw[arrow] (onnx) -- (ios)
  node[above, label, midway] {4.4\,MB ONNX};

% Below labels
\node[label, below=0.25cm of train] {mAP 94.30\%};
\node[label, below=0.25cm of export, align=center] {no-NMS\\for CoreML};
\node[label, below=0.25cm of onnx, align=center] {onnx.checker\\validated};
\node[label, below=0.25cm of ios, align=center] {React Native\\ONNX Runtime};

% Outer border
\node[draw, dashed, rounded corners, fit=(train)(export)(onnx)(ios),
      inner sep=10pt,
      label={[font=\scriptsize\itshape]above:End-to-End Deployment Pipeline}] {};
\end{tikzpicture}}
\caption{End-to-end deployment pipeline from PaddlePaddle training to
  iOS on-device inference.}
\label{fig:pipeline}
\end{figure}

Key decisions at each stage:

\begin{itemize}[leftmargin=1.5em]
  \item \textbf{Export without NMS:} The standard PicoDet export
    includes NMS in the computation graph.  For iOS deployment, we use
    the \texttt{--export\_serving\_model=False} flag to produce a
    \emph{no-NMS} graph whose outputs (raw boxes + scores) are compatible
    with CoreML's built-in NMS primitive.

  \item \textbf{ONNX Opset 11:} All operators in the no-NMS PicoDet-S
    graph map to standard ONNX Opset~11 primitives, avoiding the
    \texttt{multiclass\_nms3} custom operator that makes YOLOv3
    incompatible with CoreML.

  \item \textbf{React Native + ONNX Runtime:} The iOS app is built with
    React Native (TypeScript) and uses the
    \texttt{onnxruntime-react-native} package.  Image frames are
    preprocessed in a native Swift module
    (\texttt{ImagePreprocessor.swift}) using Metal's \texttt{vImage}
    pipeline to convert YUV frames to normalised RGB tensors on-device,
    avoiding CPU-side pixel format conversion.
\end{itemize}

\subsection{YOLOv3 Deployment Challenges}

YOLOv3 was the first model deployed; it exposed a critical
incompatibility with CoreML:

\begin{table}[H]
\centering
\caption{Comparison of deployment challenges between YOLOv3 and PicoDet-S.}
\label{tab:deploy_compare}
\begin{tabular}{lll}
\toprule
\textbf{Aspect} & \textbf{YOLOv3-MobileNetV1} & \textbf{PicoDet-S} \\
\midrule
ONNX size       & 92\,MB  & 4.4\,MB \\
Custom operator & \texttt{multiclass\_nms3} & None \\
CoreML support  & Blocked (custom op) & Full (standard ops) \\
Metal GPU accel & No (falls back to CPU) & Yes (ANE + GPU) \\
iOS FPS         & $\sim$1\,FPS & 14\,FPS (warm) \\
Input size      & 608$\times$608 & 320$\times$320 \\
\bottomrule
\end{tabular}
\end{table}

The 1\,FPS bottleneck for YOLOv3 is caused by the inability to offload
the \texttt{multiclass\_nms3} operator to Metal, forcing the entire
inference graph to run on the CPU with 32-bit float arithmetic and
608$\times$608 input tensors.

\subsection{Iterative Optimisation to 14\,FPS}

PicoDet-S deployment was improved across four engineering iterations:

\begin{figure}[H]
\centering
\begin{tikzpicture}
\begin{axis}[
  width=0.82\textwidth, height=5.5cm,
  xlabel={Iteration}, ylabel={iOS Inference FPS},
  symbolic x coords={v1,v2,v3,v4},
  xtick=data,
  xticklabels={v1 Baseline, v2 Swift, v3 CoreML, v4 Tuned},
  ymin=0, ymax=18,
  ytick={0,2,4,6,8,10,12,14,16},
  nodes near coords, nodes near coords align=above,
  nodes near coords style={font=\small\bfseries},
  ymajorgrids=true, grid style={line width=0.3pt, gray!30},
  enlarge x limits=0.3,
  title={iOS Inference Speed: Four Engineering Iterations},
  title style={font=\small\bfseries},
  tick label style={font=\scriptsize},
  label style={font=\small},
]
\addplot[
  mark=*, mark size=4pt, thick,
  color=accentGreen, fill=accentGreen,
  mark options={fill=accentGreen},
]
  coordinates {(v1,1)(v2,4)(v3,10)(v4,14)};

% Annotations
\node[font=\scriptsize, gray!70!black, align=center]
  at (axis cs:v1, 2.5) {YUV decode\\on CPU};
\node[font=\scriptsize, gray!70!black, align=center]
  at (axis cs:v2, 5.5) {Metal\\vImage};
\node[font=\scriptsize, gray!70!black, align=center]
  at (axis cs:v3, 11.5) {ONNX\\CoreML EP};
\node[font=\scriptsize, gray!70!black, align=center]
  at (axis cs:v4, 15.5) {Stable\\warm};
\end{axis}
\end{tikzpicture}
\caption{PicoDet-S inference speed on iPhone across four engineering
  iterations.  The dominant gains come from (a) offloading image
  preprocessing to the Metal vImage pipeline and (b) enabling the
  CoreML execution provider in ONNX Runtime for ANE acceleration.}
\label{fig:fps_iter}
\end{figure}

\begin{itemize}[leftmargin=1.5em]
  \item \textbf{v1 -- JavaScript baseline ($\sim$1\,FPS):}
    Frame decoding and normalisation performed in JavaScript (TypeScript),
    causing excessive memory copies across the JS bridge.

  \item \textbf{v2 -- Swift native preprocessing ($\sim$4\,FPS):}
    Moved image preprocessing to \texttt{ImagePreprocessor.swift}
    with \texttt{vImage} Metal pipeline.  Eliminates JS-bridge overhead.

  \item \textbf{v3 -- CoreML execution provider ($\sim$10\,FPS):}
    Enabled \texttt{CoreMLExecutionProvider} in ONNX Runtime, which
    compiles the ONNX graph for Apple's Neural Engine (ANE).

  \item \textbf{v4 -- Thread tuning, warm-machine stable ($\sim$14\,FPS):}
    Set inference thread count to match A-series performance core count;
    measured in warm-machine steady state (model loaded, cache warm).
\end{itemize}
